\chapter{The evidence}

This chapter reports every number the sweep produces, the exact reading beside the same corpus read on a grid, and what the comparison settles.

\section{What was swept}

Sixty four mineral names were put to the Crystallography Open Database~\cite{src:Crystallography-Open-Database} by text search. Every entry returned with right angles was fetched, tiled, and asked for its period along each of its three axes.

\begin{longtable}{@{}p{0.55\textwidth}p{0.37\textwidth}@{}}
\toprule
\endhead
names searched & 64 \\
minerals that returned a usable entry & 55 \\
structures read & 151 \\
axes measured & 453 \\
\bottomrule
\end{longtable}

One structure contributes up to three axes, and every axis is an independent question carrying its own published answer. A wrong result cannot hide in a small denominator, which a sweep buys and a sample does not. Three minerals returning three exact answers would still leave room for coincidence.

The same corpus is read twice. The gridded reading puts every site on a grid of 0.25 angstroms before the detector sees it. The exact reading carries each site as an exact integer and rounds nothing.

\section{The result}

\begin{longtable}{@{}p{0.55\textwidth}p{0.37\textwidth}@{}}
\toprule
\endhead
axes equal to the published edge & 453 of 453, 100.0 percent \\
error & zero, on every axis \\
tolerance applied & none \\
\bottomrule
\end{longtable}

Equal as integers, with no tolerance. Every coordinate is carried at $10^{-1024}$ angstroms, a period is read off the whole difference set of the axis, and the comparison against the published edge is an equality between two integers. There is no tolerance to set and no error to average, because there is no error.

The failure stated in advance is any axis whose recovered period is not equal to its published edge. A near miss counts as a miss, since nothing in this path rounds. No axis misses.

\section{The gridded reading, and what it measures}

The same 453 axes, read on the grid:

\begin{longtable}{@{}p{0.55\textwidth}p{0.37\textwidth}@{}}
\toprule
\endhead
axes inside one voxel of 0.25 angstroms & 453 of 453, 100.0 percent \\
axes equal to the published edge & 0 of 453 \\
mean absolute error & 0.0124 angstroms \\
median absolute error & 0.0108 angstroms \\
90th percentile & 0.0243 angstroms \\
99th percentile & 0.0404 angstroms \\
worst single axis & 0.0554 angstroms \\
\bottomrule
\end{longtable}

Read alone, that is a good result. Read beside the exact one, it says something else: the gridded reading never once lands on the published edge, and every figure in the table is the grid.

None of that error is in the deposit and none of it is in the detector. A cell edge is published to five decimal places and the detector compares integers. The 0.0124 is introduced in the reader, before the measurement begins. Reporting it as the accuracy of the instrument mistakes the reader for the instrument.

\section{The bias is the voxel floor}

The gridded reading carries a bias:

\begin{longtable}{@{}p{0.55\textwidth}p{0.37\textwidth}@{}}
\toprule
\endhead
mean signed error & $-0.0067$ angstroms \\
axes read long & 120 \\
axes read short & 327 \\
\bottomrule
\end{longtable}

The detector reads short about three times as often as it reads long, and the mean sits below zero by roughly half the median absolute error. A symmetric error would put those counts near even.

The voxel floor predicts a bias of about half a voxel spread over the tiling. \textbf{The test is to remove the rounding.} The exact reading of the same 453 axes has a signed error of zero, 0 axes long and 0 axes short. The bias is the voxel floor in its entirety. Nothing of it survives when sites are not assigned to cells.

That is the strongest form the test can take. Changing the rounding rule would move the bias and leave the question of how much is rule and how much is something else. Removing the quantization removes the bias completely, which settles it.

\section{The four worst gridded axes}

\begin{longtable}{@{}p{0.18\textwidth}p{0.13\textwidth}p{0.05\textwidth}p{0.15\textwidth}p{0.15\textwidth}p{0.13\textwidth}p{0.10\textwidth}@{}}
\toprule
mineral & entry & axis & published & on the grid & grid error & exact \\
\midrule
\endhead
corundum & 1544906 & c & 2.8446 & 2.9000 & $+0.0554$ & 2.8446 \\
herzenbergite & 9008295 & a & 4.1480 & 4.1000 & $-0.0480$ & 4.1480 \\
quartz & 4339381 & c & 15.4780 & 15.5250 & $+0.0470$ & 15.4780 \\
herzenbergite & 9008785 & b & 11.1800 & 11.1389 & $-0.0411$ & 11.1800 \\
\bottomrule
\end{longtable}

Herzenbergite appears twice, and it is the mineral an uncapped harmonic family reads short. It is the hardest case in the set on the grid and it is exact on both entries, as is every other axis. The exact reading has no worst axis.

The two largest grid errors have opposite signs: the worst case on the grid is not the bias but the resolution of the straddling ratio, and those are separate effects. Both are properties of a quantization the exact reading does not perform.

\section{What this table is not}

It is not a measurement of any of these minerals. Every published edge in the table was measured by somebody else, to a precision this does not reach, and is the input the recovered value is graded against. The exact reading recovers the deposited number and makes no claim beyond it: agreeing with a deposit to the digit is not the same as the deposit being right, and the bracketed uncertainty a deposit carries is a quantum fixed upstream by the people who measured it.

It is also not 453 independent structures. It is 453 axes over 151 structures over 55 minerals, and several minerals contribute many entries. Treating the 453 as independent samples of anything would overstate what a sweep of 55 minerals shows.

\section{Availability}

Each structure is a Crystallography Open Database entry, named by its entry number where the text quotes it. The detector, the exact site representation and the sweep are in the orior repository, \texttt{https://github.com/dstroy0/orior}.
